\documentclass[11pt, a4paper]{article}
\usepackage[T1]{fontenc}
\usepackage[utf8]{inputenc}
\usepackage[margin=2.5cm]{geometry}
\usepackage{microtype}
\usepackage{booktabs}
\usepackage{amsmath, amssymb}
\usepackage{mathtools}
\usepackage{enumitem}
\usepackage{parskip}
\usepackage{hyperref}
\hypersetup{colorlinks=true, linkcolor=blue, citecolor=blue, urlcolor=blue}

\title{Further Model Comparison Ideas\\[4pt]\large Internal Working Document}
\author{Hening Wang}
\date{March 2026}

\begin{document}
\maketitle

\section{Motivation}

The current model space crosses two dimensions:
\begin{center}
\begin{tabular}{lcc}
\toprule
 & Global speaker & Incremental speaker \\
\midrule
Pooled       & \checkmark & \checkmark \\
Hierarchical & \checkmark & \checkmark \\
\bottomrule
\end{tabular}
\end{center}

Both speakers share the same incremental listener, which features \emph{context-recursive} compositional semantics: the size threshold $\theta_k$ is recomputed at each composition step relative to the running posterior.
This means two mechanisms are confounded in the current incremental-vs-global comparison:

\begin{enumerate}[label=(\roman*)]
  \item \textbf{Production-level incrementality:} greedy left-to-right token selection (incremental speaker) vs.\ global utterance-level softmax (global speaker).
  \item \textbf{Listener-level context recursion:} size semantics depend on the running posterior $P_t$, making word order matter even for the global speaker.
\end{enumerate}

When the incremental hierarchical model wins LOO, we cannot tell how much of the advantage comes from (i) and how much from (ii).
The two extensions below are designed to tease these sources apart.


\section{Extension 1: Incremental Speaker with Static Semantics}

\subsection{Idea}

Remove the context recursion from the listener.
Instead of recomputing the size threshold at each composition step relative to the updated posterior $P_t$, fix the threshold once based on the \emph{original} uniform prior $P_T = P_\mathrm{prior}$.
All other aspects of the model stay the same: the incremental speaker still produces tokens left-to-right via the chain rule, and the listener still composes meanings right-to-left --- but the meanings are now \emph{static} (context-independent).

\subsection{Formal definition}

Under static semantics, the size semantic value becomes:
\[
\llbracket \text{S} \rrbracket_\mathrm{static}(s)
= \Phi\!\left(
    \frac{s_o - \theta_k(P_\mathrm{prior})}
         {w_f \sqrt{s_o^2 + \theta_k(P_\mathrm{prior})^2}}
  \right),
\]
where $\theta_k(P_\mathrm{prior})$ is computed once from the full, unfiltered context.
The listener update simplifies to a standard (non-recursive) Bayesian update:
\[
P_{t-1}(s) \propto P_t(s) \cdot \llbracket w_t \rrbracket_\mathrm{static}(s).
\]
Because $\llbracket w_t \rrbracket_\mathrm{static}$ no longer depends on $P_t$, this composition is commutative for fixed-length utterances:
\[
P(s \mid w_1, w_2) = P(s \mid w_2, w_1).
\]
The listener no longer distinguishes word order.
Any ordering preferences must therefore come entirely from the \emph{speaker} side (production-level incrementality).

\subsection{Model variants to add}

\begin{center}
\begin{tabular}{llll}
\toprule
\textbf{Model} & \textbf{Speaker} & \textbf{Semantics} & \textbf{New?} \\
\midrule
Global, static semantics       & Global       & Static & New \\
Incremental, static semantics  & Incremental  & Static & New \\
\bottomrule
\end{tabular}
\end{center}

Only the hierarchical variants are needed (they dominate pooled in all existing comparisons).

\subsection{What this test reveals}

\begin{center}
\begin{tabular}{lcc}
\toprule
\textbf{Comparison} & \textbf{Question} \\
\midrule
Inc.\ recursive vs.\ Inc.\ static & Does context recursion help the incremental speaker? \\
Global recursive vs.\ Global static & Does context recursion help the global speaker? \\
Inc.\ static vs.\ Global static & Does production incrementality help without context recursion? \\
Inc.\ recursive vs.\ Global recursive & (Already have) Combined effect \\
\bottomrule
\end{tabular}
\end{center}

Expected outcomes:
\begin{itemize}
  \item If Inc.\ recursive $\gg$ Inc.\ static: context recursion is doing heavy lifting, even within the incremental speaker.
  \item If Inc.\ static $\approx$ Global static: production-level incrementality adds nothing once you remove context recursion --- the ordering signal comes entirely from the listener.
  \item If Inc.\ static $>$ Global static: production-level incrementality contributes an independent ordering preference even with static semantics.
\end{itemize}

\subsection{Implementation}

Minimal code change: in \texttt{incremental\_semantics\_jax} and its incremental-speaker counterpart, replace the call to \texttt{compute\_size\_semantics(states, state\_prior=running\_posterior, ...)} with \texttt{compute\_size\_semantics(states, state\_prior=uniform, ...)}, fixing the prior to uniform at every composition step.
This is a one-line change per function (pass the original uniform prior instead of the current posterior).


\section{Extension 2: Lookahead Speaker ($d = 1$)}

\subsection{Idea}

Let the incremental speaker ``plan one step ahead'' before committing to a token.
Instead of choosing the token that maximises immediate informativeness, it chooses the token that leads to the best outcome after one additional step.

\subsection{Formal definition}

At step $t$, the lookahead-$d$ speaker scores candidate token $a$ using a value function:
\begin{align}
  V(u_{\leq t},\; s^*,\; 0) &= \bigl[L(s^* \mid u_{\leq t})\bigr]^\alpha, \\
  V(u_{\leq t},\; s^*,\; d) &= \max_{a'}\; V(u_{\leq t} \oplus a',\; s^*,\; d-1).
\end{align}
The local decision rule becomes:
\[
S_t(a \mid u_{<t}, s^*) \;\propto\; V(u_{<t} \oplus a,\; s^*,\; d).
\]

Three values of $d$ span the continuum:
\begin{center}
\begin{tabular}{cll}
\toprule
$d$ & Behaviour & Status \\
\midrule
0 & Greedy (current incremental speaker) & Existing \\
1 & One-step lookahead & \textbf{New} \\
2 & Full lookahead ($\approx$ global with 3 tokens) & $\approx$ Existing \\
\bottomrule
\end{tabular}
\end{center}

\subsection{Why only $d = 1$ is genuinely new}

With $T_\mathrm{max} = 3$ and vocabulary $\{S, C, F\}$:
\begin{itemize}
  \item At step 1: $d=1$ evaluates all $(a_1, a_2)$ pairs (6 paths), picks $a_1$ based on the best continuation.
  \item At step 2: $d=1$ evaluates the one remaining token --- identical to greedy.
  \item $d=2$ at step 1 evaluates all $(a_1, a_2, a_3)$ triples --- exhaustive, collapsing to global-like behaviour.
\end{itemize}

So $d=1$ differs from greedy \emph{only at step 1}.

\subsection{What this test reveals}

\begin{center}
\begin{tabular}{ll}
\toprule
\textbf{Comparison} & \textbf{Question} \\
\midrule
Inc.\ greedy vs.\ Inc.\ lookahead & Does planning ahead improve fit? \\
Inc.\ lookahead vs.\ Global & Is one-step planning closer to greedy or global? \\
\bottomrule
\end{tabular}
\end{center}

Expected outcome: if greedy $\geq$ lookahead in LOO, this supports the psychological claim that adjective ordering arises from \emph{myopic, step-by-step} pragmatic reasoning rather than forward planning.
The paper could then state:

\begin{quote}
``Even when permitted one step of lookahead planning, the speaker model does not improve its fit to human data, suggesting that ordering preferences are better explained by greedy, step-by-step pragmatic optimisation.''
\end{quote}

\subsection{Implementation}

Modify step $t=0$ in \texttt{incremental\_speaker}:
\begin{enumerate}
  \item For each candidate first token $a \in \{S, C, F\}$, simulate the listener update after $a$.
  \item For each possible second token $a'$ (given $a$), compute $[L(s^* \mid a, a')]^\alpha$.
  \item Take the max over $a'$ as the score for $a$.
  \item Steps $t \geq 1$: unchanged (greedy).
\end{enumerate}
Estimated effort: $\sim$50 lines of new code within the existing scan structure.


\section{Summary: Extended Model Space}

\begin{center}
\begin{tabular}{llll}
\toprule
\textbf{Model} & \textbf{Speaker} & \textbf{Semantics} & \textbf{Status} \\
\midrule
Global, hier.              & Global       & Recursive (current) & Existing \\
Incremental, hier.         & Incremental ($d=0$) & Recursive (current) & Existing \\
\midrule
Global, static, hier.      & Global       & Static              & \textbf{New (Ext.\ 1)} \\
Incremental, static, hier. & Incremental ($d=0$) & Static              & \textbf{New (Ext.\ 1)} \\
\midrule
Inc.\ lookahead, hier.     & Incremental ($d=1$) & Recursive (current) & \textbf{New (Ext.\ 2)} \\
\bottomrule
\end{tabular}
\end{center}

Together, these five hierarchical models allow a $2 \times 2$ factorial decomposition (speaker type $\times$ semantics type) plus a lookahead variant, enabling us to disentangle the contributions of:
\begin{enumerate}
  \item Context-recursive compositional semantics (listener side)
  \item Greedy incremental production (speaker side)
  \item Forward planning (speaker side)
\end{enumerate}

\section{Priority and Effort}

\begin{center}
\begin{tabular}{llll}
\toprule
\textbf{Extension} & \textbf{Scientific value} & \textbf{Impl.\ effort} & \textbf{Priority} \\
\midrule
Ext.\ 1 (static semantics) & High --- disentangles two confounded mechanisms & Low ($\sim$1-line change) & \textbf{1st} \\
Ext.\ 2 (lookahead $d=1$)  & Medium --- strengthens myopia claim & Medium ($\sim$50 lines) & 2nd \\
\bottomrule
\end{tabular}
\end{center}

Extension 1 should be run first: it has the highest insight-to-effort ratio and directly addresses the question of where ordering preferences come from.
Extension 2 is a nice addition once Extension 1 is done, cementing the ``greedy is best'' narrative.

\end{document}
